\section{Semisimple}

\subsection{Ideal Operations}

\begin{definition}[idealSum]
Let $I, K$ be Jordan ideals of $J$. The \textbf{sum of two ideals} is defined as:
\[ \mathsf{idealSum}(I, K) \defas \{ x \in J \mid \exists\, a \in I,\, \exists\, b \in K,\, x = a + b \} \]
This forms a Jordan ideal with:
\begin{itemize}
  \item Addition: If $x = a_1 + b_1$ and $y = a_2 + b_2$ with $a_i \in I$, $b_i \in K$, then $x + y = (a_1 + a_2) + (b_1 + b_2)$
  \item Zero: $0 = 0 + 0$ with $0 \in I$ and $0 \in K$
  \item Scalar multiplication: $r \cdot x = r \cdot a + r \cdot b$
  \item Negation: $-x = (-a) + (-b)$
  \item Jordan multiplication: $y \jmul x = (y \jmul a) + (y \jmul b)$
\end{itemize}
\end{definition}

\begin{lemma}[mem\_idealSum]
For Jordan ideals $I, K$ and $x \in J$:
\[ x \in \mathsf{idealSum}(I, K) \iff \exists\, a \in I,\, \exists\, b \in K,\, x = a + b \]
\end{lemma}

\begin{lemma}[le\_idealSum\_left]
For Jordan ideals $I, K$: $I \subseteq \mathsf{idealSum}(I, K)$.
\end{lemma}

\begin{lemma}[le\_idealSum\_right]
For Jordan ideals $I, K$: $K \subseteq \mathsf{idealSum}(I, K)$.
\end{lemma}

\begin{definition}[idealInf]
Let $I, K$ be Jordan ideals of $J$. The \textbf{intersection of two ideals} is defined as:
\[ \mathsf{idealInf}(I, K) \defas I \cap K \]
This forms a Jordan ideal.
\end{definition}

\begin{lemma}[mem\_idealInf]
For Jordan ideals $I, K$ and $x \in J$:
\[ x \in \mathsf{idealInf}(I, K) \iff x \in I \land x \in K \]
\end{lemma}

\begin{lemma}[idealInf\_le\_left]
For Jordan ideals $I, K$: $\mathsf{idealInf}(I, K) \subseteq I$.
\end{lemma}

\begin{lemma}[idealInf\_le\_right]
For Jordan ideals $I, K$: $\mathsf{idealInf}(I, K) \subseteq K$.
\end{lemma}

\begin{definition}[Independent]
Two Jordan ideals $I, K$ are \textbf{independent} if their intersection is zero:
\[ \mathsf{Independent}(I, K) \defas \mathsf{idealInf}(I, K) = \bot \]
\end{definition}

\begin{lemma}[independent\_iff]
For Jordan ideals $I, K$:
\[ \mathsf{Independent}(I, K) \iff \forall\, x,\, (x \in I \to x \in K \to x = 0) \]
\end{lemma}

\begin{definition}[IsDirectSum]
Two Jordan ideals $I, K$ form a \textbf{direct sum decomposition} if they are independent and span $J$:
\[ \mathsf{IsDirectSum}(I, K) \defas \mathsf{Independent}(I, K) \land \mathsf{idealSum}(I, K) = \top \]
\end{definition}

\begin{lemma}[isDirectSum\_iff]
For Jordan ideals $I, K$:
\[ \mathsf{IsDirectSum}(I, K) \iff \bigl(\forall\, x,\, x \in I \to x \in K \to x = 0\bigr) \land \bigl(\forall\, y \in J,\, \exists\, a \in I,\, \exists\, b \in K,\, y = a + b\bigr) \]
\end{lemma}

\subsection{Semisimple Jordan Algebras}

\begin{definition}[IsSemisimpleJordan]
A Jordan algebra $J$ is \textbf{semisimple} if:
\begin{itemize}
  \item (Nontrivial) $\exists\, a \in J,\, a \ne 0$
  \item (Complement property) For every Jordan ideal $I$ of $J$, there exists a Jordan ideal $K$ such that $\mathsf{IsDirectSum}(I, K)$
\end{itemize}
\end{definition}

\begin{proposition}
Every semisimple Jordan algebra $J$ is nontrivial (i.e., $\exists\, a, b \in J$ with $a \ne b$).
\end{proposition}

\begin{theorem}[jone\_ne\_zero]
In a semisimple Jordan algebra $J$, the identity is nonzero: $\jone \ne 0$.
\end{theorem}
\begin{proof}
Suppose $\jone = 0$. By semisimplicity, there exists $a \in J$ with $a \ne 0$. But then:
\[ a = a \jmul \jone = a \jmul 0 = 0 \]
which contradicts $a \ne 0$.
\end{proof}

\begin{theorem}[IsSimpleJordan.isSemisimpleJordan]
Every simple Jordan algebra is semisimple.
\end{theorem}
\begin{proof}
Let $J$ be a simple Jordan algebra and $I$ a Jordan ideal. By simplicity, either $I = \bot$ or $I = \top$.
\begin{itemize}
  \item If $I = \bot$: Take $K = \top$. Then $\mathsf{idealInf}(\bot, \top) = \bot$ and $\mathsf{idealSum}(\bot, \top) = \top$.
  \item If $I = \top$: Take $K = \bot$. Then $\mathsf{idealInf}(\top, \bot) = \bot$ and $\mathsf{idealSum}(\top, \bot) = \top$.
\end{itemize}
In both cases, $\mathsf{IsDirectSum}(I, K)$ holds.
\end{proof}

\begin{theorem}[IsSemisimpleJordan.mk']
Let $J$ be a nontrivial Jordan algebra. If every Jordan ideal $I$ of $J$ has a complement $K$ such that $\mathsf{IsDirectSum}(I, K)$, then $J$ is semisimple.
\end{theorem}
